% GENERATED FILE. Edit canonical lessons, then run npm run generate:books.
% canonical-source-hash: fnv1a64:b48c0a731caa13c9
% canonical-lessons: BN-C238-mora, BN-C238-barano, BN-C238-komano, BN-C238-tara-kora, BN-C238-bishvas-kora

\chapter{To die, To increase, To reduce, To chase, To believe}
\label{ch:bn-to-die-to-increase-to-reduce-to-chase-to-believe}
\hlchaptermodality{238}

\begin{chapteropening}
\textbf{By the end of this chapter:} \emph{I can say to die, to increase, to reduce, to chase and to believe.}

Complete the last lesson of chapter 238: I can say to die, to increase, to reduce, to chase and to believe.
\end{chapteropening}

\section[\texorpdfstring{\bn{মরা} (môrā)}{মরা (môrā)}]{\bn{মরা} (môrā) — to die}

\label{lesson:BN-C238-mora}



\begin{quote}
Before the new one: say the Bengali for to cross, then the Bengali for to go on an outing.
\end{quote}

\subsection*{You'll want to know: \bn{মরা}}
\textbf{\bn{মরা}} — \emph{môrā} — \textquotedblleft{}to die\textquotedblright{}.

\begin{cousinweb}
\textbf{\bn{মরা}} (\emph{môrā}) is a tadbhava, a word inherited from Sanskrit and reshaped by sound change. It comes from \emph{marate}, \textquotedblleft{}will die\textquotedblright{} (CDIAL 9871). CDIAL lists Bengali \emph{marā}, \textquotedblleft{}dies\textquotedblright{}, and thinks it possible or even probable that the verb continues this form directly.
\end{cousinweb}

\begin{grammarlens}[title={What you've built}]
One more word for this chapter.
\end{grammarlens}

\subsection*{Your turn}
\begin{itemize}
\raggedright
  \item \emph{Say it:} \emph{môrā}
  \item \emph{Say it:} \emph{môrā}, once more
  \item \emph{Say it:} say \emph{môrā}
  \item \emph{Recall:} say the Bengali for to cross, then the Bengali for to go on an outing, then say \emph{môrā} again
\end{itemize}

\begin{tcolorbox}[breakable,colback=teal!4,colframe=teal!35!black,title={Before you move on}]
What does \bn{মরা} mean? (\textquotedblleft{}To die\textquotedblright{}.) Say it once more.
\end{tcolorbox}
\section[\texorpdfstring{\bn{বাড়ানো} (bāṛāno)}{বাড়ানো (bāṛāno)}]{\bn{বাড়ানো} (bāṛāno) — to increase}

\label{lesson:BN-C238-barano}



\begin{quote}
Before the new one: say the Bengali for to go on an outing, then the Bengali for to die.
\end{quote}

\subsection*{You'll want to know: \bn{বাড়ানো}}
\textbf{\bn{বাড়ানো}} — \emph{bāṛāno} — \textquotedblleft{}to increase\textquotedblright{}.

\begin{cousinweb}
\textbf{\bn{বাড়ানো}} (\emph{bāṛāno}) is a tadbhava, a word inherited from Sanskrit and reshaped by sound change. It comes from \emph{vardhate}, \textquotedblleft{}grows, increases\textquotedblright{} (CDIAL 11376). CDIAL lists Bengali \emph{bāṛā}, \textquotedblleft{}to increase\textquotedblright{}; \emph{bāṛāno} is its causing form.
\end{cousinweb}

\begin{grammarlens}[title={What you've built}]
One more word for this chapter.
\end{grammarlens}

\subsection*{Your turn}
\begin{itemize}
\raggedright
  \item \emph{Say it:} \emph{bāṛāno}
  \item \emph{Say it:} \emph{bāṛāno}, once more
  \item \emph{Say it:} say \emph{bāṛāno}
  \item \emph{Recall:} say the Bengali for to go on an outing, then the Bengali for to die, then say \emph{bāṛāno} again
\end{itemize}

\begin{tcolorbox}[breakable,colback=teal!4,colframe=teal!35!black,title={Before you move on}]
What does \bn{বাড়ানো} mean? (\textquotedblleft{}To increase\textquotedblright{}.) Say it once more.
\end{tcolorbox}
\section[\texorpdfstring{\bn{কমানো} (kômāno)}{কমানো (kômāno)}]{\bn{কমানো} (kômāno) — to reduce}

\label{lesson:BN-C238-komano}



\begin{quote}
Before the new one: say the Bengali for to die, then the Bengali for to increase.
\end{quote}

\subsection*{You'll want to know: \bn{কমানো}}
\textbf{\bn{কমানো}} — \emph{kômāno} — \textquotedblleft{}to reduce\textquotedblright{}.

\begin{cousinweb}
\textbf{\bn{কমানো}} (\emph{kômāno}) is built on \emph{kôm}, \textquotedblleft{}less\textquotedblright{}, usually traced to Persian \emph{kam}, \textquotedblleft{}little, less\textquotedblright{}, with the causing ending \emph{-āno}: \textquotedblleft{}to make less\textquotedblright{}.
\end{cousinweb}

\begin{grammarlens}[title={What you've built}]
One more word for this chapter.
\end{grammarlens}

\subsection*{Your turn}
\begin{itemize}
\raggedright
  \item \emph{Say it:} \emph{kômāno}
  \item \emph{Say it:} \emph{kômāno}, once more
  \item \emph{Say it:} say \emph{kômāno}
  \item \emph{Recall:} say the Bengali for to die, then the Bengali for to increase, then say \emph{kômāno} again
\end{itemize}

\begin{tcolorbox}[breakable,colback=teal!4,colframe=teal!35!black,title={Before you move on}]
What does \bn{কমানো} mean? (\textquotedblleft{}To reduce\textquotedblright{}.) Say it once more.
\end{tcolorbox}
\section[\texorpdfstring{\bn{তাড়া} \bn{করা} (tāṛā kôrā)}{তাড়া করা (tāṛā kôrā)}]{\bn{তাড়া} \bn{করা} (tāṛā kôrā) — to chase}

\label{lesson:BN-C238-tara-kora}



\begin{quote}
Before the new one: say the Bengali for to increase, then the Bengali for to reduce.
\end{quote}

\subsection*{You'll want to know: \bn{তাড়া} \bn{করা}}
\textbf{\bn{তাড়া} \bn{করা}} — \emph{tāṛā kôrā} — \textquotedblleft{}to chase\textquotedblright{}.

\begin{cousinweb}
\textbf{\bn{তাড়া} \bn{করা}} (\emph{tāṛā kôrā}) is \emph{tāṛā}, \textquotedblleft{}a chase\textquotedblright{}, which belongs with Bengali \emph{tāṛāna}, \textquotedblleft{}to drive away\textquotedblright{}, from Sanskrit \emph{tāḍayati}, \textquotedblleft{}beats, strikes\textquotedblright{} (CDIAL 5752), + \emph{kôrā}, \textquotedblleft{}to do\textquotedblright{}, from Sanskrit \emph{karoti}, \textquotedblleft{}does\textquotedblright{} (CDIAL 2814).
\end{cousinweb}

\begin{grammarlens}[title={What you've built}]
One more word for this chapter.
\end{grammarlens}

\subsection*{Your turn}
\begin{itemize}
\raggedright
  \item \emph{Say it:} \emph{tāṛā kôrā}
  \item \emph{Say it:} \emph{tāṛā kôrā}, once more
  \item \emph{Say it:} say \emph{tāṛā kôrā}
  \item \emph{Recall:} say the Bengali for to increase, then the Bengali for to reduce, then say \emph{tāṛā kôrā} again
\end{itemize}

\begin{tcolorbox}[breakable,colback=teal!4,colframe=teal!35!black,title={Before you move on}]
What does \bn{তাড়া} \bn{করা} mean? (\textquotedblleft{}To chase\textquotedblright{}.) Say it once more.
\end{tcolorbox}
\section[\texorpdfstring{\bn{বিশ্বাস} \bn{করা} (bishvās kôrā)}{বিশ্বাস করা (bishvās kôrā)}]{\bn{বিশ্বাস} \bn{করা} (bishvās kôrā) — to believe, to trust}

\label{lesson:BN-C238-bishvas-kora}



\begin{quote}
Before the new one: say the Bengali for to reduce, then the Bengali for to chase.
\end{quote}

\subsection*{You'll want to know: \bn{বিশ্বাস} \bn{করা}}
\textbf{\bn{বিশ্বাস} \bn{করা}} — \emph{bishvās kôrā} — \textquotedblleft{}to believe, to trust\textquotedblright{}.

\begin{cousinweb}
\textbf{\bn{বিশ্বাস} \bn{করা}} (\emph{bishvās kôrā}) is \emph{bishvās}, \textquotedblleft{}trust\textquotedblright{}, Sanskrit \emph{viśvāsa} taken over whole, + \emph{kôrā}, \textquotedblleft{}to do\textquotedblright{}, from Sanskrit \emph{karoti}, \textquotedblleft{}does\textquotedblright{} (CDIAL 2814).
\end{cousinweb}

\begin{grammarlens}[title={What you've built}]
One more word for this chapter.
\end{grammarlens}

\subsection*{Your turn}
\begin{itemize}
\raggedright
  \item \emph{Say it:} \emph{bishvās kôrā}
  \item \emph{Say it:} \emph{bishvās kôrā}, once more
  \item \emph{Say it:} say \emph{bishvās kôrā}
  \item \emph{Recall:} say the Bengali for to reduce, then the Bengali for to chase, then say \emph{bishvās kôrā} again
\end{itemize}

\begin{tcolorbox}[breakable,colback=teal!4,colframe=teal!35!black,title={Before you move on}]
What does \bn{বিশ্বাস} \bn{করা} mean? (\textquotedblleft{}To believe, to trust\textquotedblright{}.) Say it once more.
\end{tcolorbox}
